\documentclass{article}
\usepackage{amsmath}
\usepackage{amssymb}
\usepackage[margin=1in]{geometry}

\title{Counterexample Note: Interior Feasibility Is Not Implied by Local Positivity}
\author{}
\date{\today}

\begin{document}
\maketitle

\section*{What we attempted}
Before proving the \(m \le M_k(N)\) lemma, we tested whether the converse could be made unconditional by dropping the interior-feasibility assumption.

The hoped-for claim was:
\[
\mathcal{F}_\Phi \neq \varnothing
\quad\Longrightarrow\quad
\exists P^\circ \in \mathcal{F}_\Phi \text{ with } \operatorname{supp}(P^\circ)=T_\Phi,
\]
where
\[
T_\Phi := \{y : \forall i,\ \phi_i(y_{S_i})>0\}.
\]
If true, this would remove the extra regularity assumption from the maximal-support step.

\section*{Counterexample}
Take \(n=3\), \(k=2\), and constraints \(\Phi=\{\phi_{XY},\phi_{XZ},\phi_{YZ}\}\) with pairwise marginals
\[
\phi_{XY}=
\begin{array}{c|cc}
 & Y=0 & Y=1\\\hline
X=0 & \tfrac{9}{20} & \tfrac{1}{20}\\
X=1 & \tfrac{1}{20} & \tfrac{9}{20}
\end{array},
\quad
\phi_{XZ}=
\begin{array}{c|cc}
 & Z=0 & Z=1\\\hline
X=0 & \tfrac{2}{5} & \tfrac{1}{10}\\
X=1 & \tfrac{1}{10} & \tfrac{2}{5}
\end{array},
\]
\[
\phi_{YZ}=
\begin{array}{c|cc}
 & Z=0 & Z=1\\\hline
Y=0 & \tfrac{7}{20} & \tfrac{3}{20}\\
Y=1 & \tfrac{3}{20} & \tfrac{7}{20}
\end{array}.
\]
Every local entry is strictly positive, so \(T_\Phi=\{0,1\}^3\).

Let \(p_{xyz}=\Pr(X=x,Y=y,Z=z)\). Solving the linear marginal equations gives
\[
\begin{aligned}
p_{000}&=\tfrac{7}{10}-t, & p_{001}&=-\tfrac14+t, & p_{010}&=-\tfrac{3}{10}+t, & p_{011}&=\tfrac{7}{20}-t,\\
p_{100}&=-\tfrac{7}{20}+t, & p_{101}&=\tfrac25-t, & p_{110}&=\tfrac{9}{20}-t, & p_{111}&=t.
\end{aligned}
\]
Nonnegativity imposes \(t\ge \tfrac{7}{20}\) (from \(p_{100}\ge0\)) and \(t\le \tfrac{7}{20}\) (from \(p_{011}\ge0\)), hence \(t=\tfrac{7}{20}\) is forced.

Therefore the unique feasible joint is
\[
(p_{000},p_{001},p_{010},p_{011},p_{100},p_{101},p_{110},p_{111})
=\Big(\tfrac{7}{20},\tfrac{1}{10},\tfrac{1}{20},0,0,\tfrac{1}{20},\tfrac{1}{10},\tfrac{7}{20}\Big),
\]
which has zeros at \(011\) and \(100\).

So \(\mathcal{F}_\Phi\neq\varnothing\), but no feasible \(P^\circ\) is strictly positive on \(T_\Phi\). Interior feasibility fails.

\section*{Why this is a dead end}
This shows that \emph{local positivity of marginals} does not guarantee a feasible distribution with full support on \(T_\Phi\). Global compatibility can force the feasible set onto a boundary face of the simplex, even when every locally checked tuple is allowed.

So the route ``consistency \(+\) local positivity \(\Rightarrow\) maximal support'' is false in general, and the converse argument needs an extra hypothesis (or a stronger structural condition) beyond consistency.

\section*{Python code used}
The script below (standard library only) computes the affine solution family and the feasible interval for \(t\) exactly over \(\mathbb{Q}\).

\begin{verbatim}
from fractions import Fraction
from itertools import product

states=list(product([0,1],[0,1],[0,1]))
idx={s:i for i,s in enumerate(states)}

mxy={(0,0):Fraction(9,20),(0,1):Fraction(1,20),(1,0):Fraction(1,20),(1,1):Fraction(9,20)}
mxz={(0,0):Fraction(2,5),(0,1):Fraction(1,10),(1,0):Fraction(1,10),(1,1):Fraction(2,5)}
myz={(0,0):Fraction(7,20),(0,1):Fraction(3,20),(1,0):Fraction(3,20),(1,1):Fraction(7,20)}

A=[]; b=[]
A.append([Fraction(1) for _ in states]); b.append(Fraction(1))
for uv in product([0,1],[0,1]):
    row=[Fraction(0) for _ in states]
    for s,i in idx.items():
        if (s[0],s[1])==uv: row[i]=Fraction(1)
    A.append(row); b.append(mxy[uv])
for uv in product([0,1],[0,1]):
    row=[Fraction(0) for _ in states]
    for s,i in idx.items():
        if (s[0],s[2])==uv: row[i]=Fraction(1)
    A.append(row); b.append(mxz[uv])
for uv in product([0,1],[0,1]):
    row=[Fraction(0) for _ in states]
    for s,i in idx.items():
        if (s[1],s[2])==uv: row[i]=Fraction(1)
    A.append(row); b.append(myz[uv])

m=len(A); n=len(states)
M=[A[i]+[b[i]] for i in range(m)]

# RREF over Fractions
r=0
pivot=[]
for c in range(n):
    piv=None
    for i in range(r,m):
        if M[i][c] != 0:
            piv=i; break
    if piv is None:
        continue
    M[r],M[piv]=M[piv],M[r]
    pv=M[r][c]
    M[r]=[x/pv for x in M[r]]
    for i in range(m):
        if i!=r and M[i][c]!=0:
            fac=M[i][c]
            M[i]=[M[i][j]-fac*M[r][j] for j in range(n+1)]
    pivot.append(c)
    r+=1

free=[j for j in range(n) if j not in pivot]
j=free[0]  # here, p_111 = t
print("Free variable =", states[j])
for i,c in enumerate(pivot):
    coeff = -M[i][j]
    print(f"p{states[c]} = {M[i][n]} + ({coeff})*t")

# Interval from nonnegativity A0 + B t >= 0
low=None; high=None
for i in range(n):
    if i==j:
        A0=Fraction(0); B=Fraction(1)
    else:
        row_i = pivot.index(i)
        A0=M[row_i][n]
        B=-M[row_i][j]
    if B==0:
        assert A0 >= 0
    elif B>0:
        bound = -A0/B
        low = bound if low is None or bound>low else low
    else:
        bound = -A0/B
        high = bound if high is None or bound<high else high

print("t in [", low, ",", high, "]")
\end{verbatim}

Expected final line:
\begin{verbatim}
t in [ 7/20 , 7/20 ]
\end{verbatim}

\end{document}
